%% notes-tex/028-knockout-expression-metabolic/sections/abstract.tex
%% SOURCE: experiments/028-knockout-expression/results/expression_metabolic_yko_correlation.json
A deletion's gene-level readout carries its metabolic phenotype, and how much of it
depends on which readout. On the deletions the panels share, a ridge from the Messner
2023 proteome to the 19 Mulleder 2016 amino-acid pools reads a median out-of-fold
Pearson of $0.41$ over 4{,}400 deletions, every amino acid above its permutation null,
proline at $0.61$ through the prolyl-tRNA synthetase and alanine at $0.55$ through
\gene{ALT1}; the same ridge from the Kemmeren 2014 transcriptome reads $0.18$ over
1{,}416 deletions, and from the Nadal-Ribelles 2025 perturb-seq pseudobulk $0.10$ over
2{,}113. The single features behind the strong reads are the pathway's own enzymes and
synthetases, with the sign a pool-and-flux picture predicts. On the betaxanthin titer
proxy of Cachera 2023 the three sources read alike, $0.43$, $0.35$ and $0.41$, on a
respiratory signature; the beta-carotene colony score reads $0.07$ to $0.16$. Alike
transcriptomes have alike amino-acid profiles (Mantel Spearman $-0.20$ over a million
Kemmeren pairs against a null of $0.03$), while the proteome, which predicts the
individual amino acids best, has the weakest whole-profile read ($-0.05$).

Nadal-Ribelles does not replicate itself across batches and barely agrees with either
other source on the same deletions, yet it reads the betaxanthin phenotype as well as
they do and half of the transcriptome's amino-acid signal. Agreement with an earlier
compendium is therefore not the test of whether a readout is worth training on; the
metabolic read is an independent one, and on it the perturb-seq atlas carries a
large-effect program and little of the smaller ones. Zelezniak 2018 and da Silveira
2014 share under 125 deletions with any source and read null, as their size predicts.
Two things are not measured: whether a shared plate structure between the Mulleder and
Messner panels contributes to the proteome's read, and whether the encoder trained on
the proteome carries the amino acids in its cell vector.
